%%
%% fall-lecture19.tex
%% 
%% Made by Alex Nelson <pqnelson@gmail.com>
%% Login   <alex@lisp>
%% 
%% Started on  2026-11-10T11:44:10-0800
%% Last update 2026-11-10T11:44:10-0800
%% 

\lecture[Carath\'{e}odory]{}

\begin{node}
We stated ``the'' Carath\'{e}odory theorem last time. It's the main
procedure we use to construct measures. Let us restate it:
\end{node}

\begin{theorem}
Let $\mathcal{R}$ be an algebra, let $\lambda$ be a pre-measure on
$\mathcal{R}$ [so $\lambda\colon\mathcal{R}\to[0,\infty]$ and
  $\lambda$ is $\sigma$-additive when sensible: if $E=\bigsqcup E_{j}$ is the
  countable union of disjoint $E_{j}\in\mathcal{R}$ and if
  $E\in\mathcal{R}$, then $\lambda(E)\leq\sum_{j}\lambda(E_{j})$].
Then:
\begin{enumerate}
\item $\lambda$ extends to a measure on $\sigma(\mathcal{R})$, and
\item if $\lambda$ is $\sigma$-finite, then this extension (to a
  measure) is unique.
\end{enumerate}
\end{theorem}

\begin{definition}
Let $X$ be a set. An \define{Outer Measure} of $X$ is a function
$\nu\colon\powerset{X}\to[0,\infty]$ defined on the powerset of $X$
such that
\begin{enumerate}
\item $\nu(\emptyset)=0$
\item \textsc{Monotonicity}: if $A\subset B\subset X$, then $\nu(A)\leq\nu(B)$
\item \textsc{Subadditivity}: for any countable family $A_{j}$ of
  subsets of $X$, we have $\nu(\bigsqcup A_{j})\leq\sum_{j}\nu(A_{j})$.
\end{enumerate}
\end{definition}

\begin{remark}
\begin{enumerate}
\item An outer measure is defined on \emph{all} subsets of $X$, not on
  a particular algebra of $X$ (nor on a particular $\sigma$-algebra of
  $X$) but \emph{all} $X$.
\item Some authors (like Folland) require $X$ be nonempty. This is not
  necessary, but the $X=\emptyset$ case is really boring because
  $\powerset{X}=\{\emptyset\}$ so $\nu(A)=0$ for all $A\subset X$ in
  this situation.
\item Usually, we obtain outer measures by starting with a family
  $\mathcal{E}$ of ``elementary sets'' on which outer measures are
  defined (like rectangles in the plane), then approximate arbitrary
  sets ``from the outside'' by countable unions of members of
  $\mathcal{E}$. 
\end{enumerate}
\end{remark}

\begin{example}
Let $\mathcal{R}$ be an algebra of $X$.
The only outer measure we will ever need is the \define{Covering Outer Measure}
$\nu=\lambda^{*}$. Then for any set $A\subset X$ we have
\begin{equation}
\nu(A)=\inf\{\sum\lambda(E_{j})\mid A\subset\bigcup_{j}E_{j}, E_{j}\in\mathcal{R}\}.
\end{equation}
So we consider the covering of $A$ by elements $E_{j}$ from some
algebra, then compute $\sum\lambda(E_{j})$ and take the infinimum of
all this value computed over all such coverings.

We claim that $\nu$ really is an outer measure. We need to show it
satisfies subadditivity. Let $A_{j}$ be an arbitrary family of subsets
of $X$. Then we simply take $E_{i,j}\in\mathcal{R}$ which cover each
$A_{j}\subset\bigcup_{i}E_{i,j}$, and the result follows immediately.
\end{example}

\begin{definition}[Carath\'{e}odory condition]
We start with an outer measure $\mu$ (the textbook writes it as
$\mu^{*}$, but if we use that starred notation\dots things will get
too starry). Then by definition, the set $A\subset X$ is
\define{$\mu$-Measurable} if
\begin{equation}
\forall E\subset X\ldotp\mu(E)=\mu(E\cap A)+\mu(E\cap A^{\complement}).
\end{equation}
In other words, $A$ splits $E$ into ``inner'' and ``outer'' parts.
(Folland requires $\forall E\in\mathcal{R}\ldotp(\dots)$ for this
condition, but I do not think we need such caution?)
\end{definition}

\begin{definition}
We say $\mu$ is \define{Complete} if its $\sigma$-algebra contains all
subsets of all null sets.
\end{definition}

\begin{theorem}[Carath\'{e}odory's famous theorem]
\begin{enumerate}
\item The collection of $\mu$-measurable sets
  \begin{equation}
\mathcal{M}=\{A\subset X\mid A\mbox{ is $\mu$-measurable}\}
  \end{equation}
  is a $\sigma$-algebra.
\item The restriction of $\mu$ to $\mathcal{M}$ is a complete measure.
\end{enumerate}
\end{theorem}

\begin{proof}
This requires 5 steps.

\textsc{Step 1: $\mathcal{M}$ is an algebra.}
Suffices to show: Let $A,B\in\mathcal{M}$, then $A\cup B\in\mathcal{M}$.

Proof: $\forall E\subset X\ldotp\mu(E\cap A)+\mu(E\cap A^{\complement})=\mu(E)$
since $A\in\mathcal{A}$. Then
\begin{equation}\label{eq:math110a:lec19:eq1}
\mu(E)=\mu((E\cap A)\cap B)+\mu((E\cap A)\cap B^{\complement})
+ \mu((E\cap A^{\complement})\cap B)+\mu((E\cap A^{\complement})\cap B^{\complement})
\end{equation}
Observe $E\cap(A\cup B)=(E\cap A\cap B)\cup(E\cap A\cap B^{\complement})\cup(E\cup A^{\complement}\cup B)$
since $A\cup B=(A\cap B)\cup(A^{\complement}\cap B)\cup(A\cap B^{\complement})$.
Then by subadditivity
\begin{equation}\label{eq:math110a:lec19:eq2}
\mu(E\cap(A\cup B))\leq\mu(E\cap A\cap B)+\mu(E\cap
A^{\complement}\cap B) + \mu(E\cap A\cap B^{\complement})
\end{equation}
Combining Eq~\eqref{eq:math110a:lec19:eq1} and Eq~\eqref{eq:math110a:lec19:eq2}
we get
\begin{equation}
\mu(E)\geq\mu(E\cap A^{\complement}\cap
B^{\complement})+\mu(E\cap(A\cup B))=\mu(E\cap(A\cup B)^{\complement})+\mu(E\cap(A\cup B))
\end{equation}
so $A\cup B$ splits $E$ for arbitrary $E$. This means $A\cup B\in\mathcal{M}$.

\textsc{Step 2: $\mu$ is finitely additive on $\mathcal{M}$.}
Let $A,B\in\mathcal{M}$. Assume they are disjoint $A\cap B=\emptyset$.
We will show $\mu(A\cup B)=\mu(A)+\mu(B)$ (which suffices to prove
finite additivity).

We see
\begin{subequations}
  \begin{align}
\mu(A\cup B)
&=\mu(A\cap(A\cup B))+\mu(A^{\complement}\cap(A\cup
B))\quad\mbox{since $A$ splits $A\cup B$}\\
&=\mu(A)+\mu(B)\quad\mbox{since $A\cap B=\emptyset$}.
  \end{align}
\end{subequations}
Then
\begin{equation}
A^{\complement}\cap(A\cup B)=(A^{\complement}\cap
A)\cup(A^{\complement}\cap B)=\emptyset\cup B=B.
\end{equation}
Hence $\mu$ is finite additive on $\mathcal{M}$.

\textsc{Step 3: $\mathcal{M}$ is a $\sigma$-Algebra.}
Consider disjoint $A_{j}\in\mathcal{M}$ and write $B_{n}=\bigsqcup^{n}_{j=1}A_{j}$
and $B=\bigsqcup^{\infty}_{j=1}A_{j}$. Let $E\subset X$ be arbitrary.
We have
\begin{subequations}
  \begin{align}
\mu(E\cap B_{n}) &= \mu(E\cap B_{n}\cap A_{n})+\mu(E\cap B_{n}\cap A_{n}^{\complement})\\
&=\mu(E\cap A_{n})+\mu(A\cap B_{n-1}),
  \end{align}
\end{subequations}
so we have a telescopic situation. By induction,
\begin{equation}
\mu(E\cap B_{2})=\mu(E\cap A_{2})+\mu(E\cap A_{1}),
\end{equation}
then
\begin{equation}
\mu(E\cap B_{n})=\sum^{n}_{j=1}\mu(E\cap A_{j})
\end{equation}
by induction, and therefore
\begin{subequations}
  \begin{align}
\mu(E) &= \mu(E\cap B_{n})+\mu(E\cap B_{n}^{\complement})\\
&\geq[\sum^{n}_{j=1}\mu(E\cap A_{j})]+\mu(E\cap B^{\complement})
  \end{align}
\end{subequations}
replacing $B_{n}$ by $B$ in the second term and noting
$B^{\complement}\subset B_{n}^{\complement}$. Then we take the limit
$n\to\infty$ to give us
\begin{subequations}
  \begin{align}
\mu(E)&\geq[\sum^{\infty}_{j=1}\mu(E\cap A_{j})]+\mu(E\cap B^{\complement})\\
&\geq\left(\cup^{\infty}_{j=1}(E\cap A_{j})\right)+\mu(E\cap B^{\complement})\quad\mbox{by subadditivity}\\
&\quad=\mu(E\cap B)+\mu(E\cap B^{\complement})\quad\mbox{by def of $B$}\\
&\quad\phantom{=\mu(E\cap B)}\quad\geq\mu(E)\quad\mbox{by subadditivity}.
  \end{align}
\end{subequations}
Since we have several inequalities of the form $\mu(E)\geq\dots\geq\mu(E)$,
they must be equalities. Therefore $\mu$ is defined on all sets and
$\mu$ is subadditivity. Then we have
\begin{equation}
\mu(E\cap B)+\mu(E\cap B^{\complement})=\mu(E),
\end{equation}
which implies $B$ is measurable. In particular, $B\in\mathcal{M}$, so
$\mathcal{M}$ is $\sigma$-algebra.

\textsc{Step 4: $\mu$ restricted to $\mathcal{M}$ is countable additive.}
Let $E_{j}\in\mathcal{M}$ be a sequence of pairwise disjoint sets, and
$E=\bigcup_{j}E_{j}$. Then for any $n\in\NN$, we have:
\begin{subequations}
  \begin{align}
\mu|_{\mathcal{M}}(E) &= \mu(E)\\
&\geq\mu|_{\mathcal{M}}(E)\left(\bigcup^{n}_{j=1}E_{j}\right)\\
&=\mu\left(\bigcup^{n}_{j=1}E_{j}\right)\\
&=\sum^{n}_{j=1}\mu(E_{j}).
  \end{align}
\end{subequations}
Taking the limit $n\to\infty$, we have
\begin{equation}
\mu(E)\geq\sum^{\infty}_{j=1}\mu(E_{j}).
\end{equation}
We also see that
\begin{equation}
\mu(E)=\mu|_{\mathcal{M}}(E)\leq\sum^{\infty}_{j=1}\mu|_{\mathcal{M}}(E_{j})=\sum^{\infty}_{j=1}\mu(E_{j}).
\end{equation}
This means $\mu(E)=\sum^{\infty}_{j=1}\mu(E_{j})$, hence $\mu$ is
countably additive.

\textsc{Step 5: $\mu$ restricted to $\mathcal{M}$ is complete.}
We see $\mu|_{\mathcal{M}}(\emptyset)=0$ and for all $E\in\mathcal{M}$
we see $\mu|_{\mathcal{M}}(E)\geq0$, so $\mu|_{\mathcal{M}}$ is a
measure on $\mathcal{M}$ (huzzah).

Now to show it is complete: Let $E$ be a null set $\mu|_{\mathcal{M}}(E)=0$. Let
$F\subset E$. Then for any $A\subset X$, we have
\begin{subequations}
  \begin{align}
\mu(A) &\geq\mu(A\cap F^{\complement})\\
&=\mu(A\cap F^{\complement})+\mu(A\cap E)\\
&\geq\mu(A\cap F^{\complement})+\mu(A\cap F),
  \end{align}
\end{subequations}
so $F\in\mathcal{M}$. Since $F$ was an arbitrary subset of an
arbitrary null set $E$, this means that $\mu|_{\mathcal{M}}$ is a
complete measure.
\end{proof}